% =====================================================================
\BookChapter[Java]{Java: Write Once, Run Anywhere}
\label{ch:java}

\begin{researchquestion}
How can the same compiled program run on any computer, and where does an engineer meet such programs?
\end{researchquestion}
\OrientationNote
% =====================================================================

\begin{factsheet}{Java}
First appeared & 1995; long-term-support releases Java 8, 11, 17, 21, 25\\
Designed by & James Gosling and team, Sun Microsystems (now owned by Oracle)\\
Paradigms & object-oriented, generic, functional elements (lambdas, streams)\\
Typing & static, strong\\
Execution & compiled to bytecode, run and JIT-compiled by the Java Virtual Machine (JVM)\\
File extensions & \file{.java} (source), \file{.class} (bytecode), \file{.jar} (archive)\\
Implementations & OpenJDK (free) and builds of it: Eclipse Temurin, Oracle, Amazon Corretto, \ldots\\
Used for & enterprise and banking systems, Android apps, large web back ends, big-data
  tools, scientific desktop tools, smart cards, education\\
Famous software & Android apps, Minecraft, Apache Hadoop/Spark/Kafka, ImageJ (image
  analysis), MATSim (traffic simulation), JMol, IntelliJ IDEA, Ghidra\\
\end{factsheet}

Java is one of the most widely used programming languages in the world. Its promise,
``write once, run anywhere'', means that the same compiled program runs on any computer
with a Java Virtual Machine. Java dominates large business software and Android apps. In
engineering and science it is less common than \py, C++ or MATLAB, but you will meet it in
data-processing infrastructure, in some simulation and image-analysis tools, and in the
first-year programming courses of many universities.

\LearningOutcomes{Compile and run a Java program on the Java Virtual Machine; read Java classes, interfaces and records; follow a calculation written with streams; and say where Java is used in engineering and where it is not.}

\section{History and philosophy}

\begin{history}
In 1991 a small team at Sun Microsystems led by James Gosling started a project to
program consumer electronics such as set-top boxes. Because these devices used many
different processors, they designed a language that compiled to bytecode for a virtual
machine rather than to the machine code of one processor. The language, first called Oak,
was renamed Java (after the coffee) and released in 1995, just as the World Wide Web took
off: Java ``applets'' could run inside web pages on any computer. Applets are long gone,
but the virtual-machine idea conquered servers, and from 2008 Java became the main language
of Android phones.
\end{history}

Java was designed as a simpler and safer C++: no pointers that can be manipulated, no
manual memory management (a garbage collector frees unused objects), array bounds always
checked, and no multiple inheritance of classes. Everything is organised in classes. The
language changed slowly for many years, but since 2017 new versions appear every six
months, and modern Java has records, type inference (\code{var}), lambdas, pattern
matching and text blocks.

\section{Getting started}

Install a Java Development Kit (JDK), for example Eclipse Temurin from
\url{https://adoptium.net}. IntelliJ IDEA, Eclipse and VS Code are common editors.
\codefile{Java}{java/Hello.java}{hello world}
\begin{lstlisting}[style=shell]
$ javac Hello.java      # compile: creates Hello.class (bytecode)
$ java Hello            # run it on the JVM
$ java Hello.java       # Java 11+: compile and run a single file in one step
\end{lstlisting}
\outfile{java/Hello.out}
Every Java program needs a class (here \code{Hello}, which must be in a file called
\file{Hello.java}) with a method \code{public static void main(String[] args)}. Each word
has a meaning: \code{public} (callable from outside), \code{static} (belongs to the class,
no object needed), \code{void} (returns nothing), and \code{String[] args} (the
command-line arguments). Recent Java versions (Java 25) also allow a much shorter form for
small programs, \code{void main() \{ IO.println("Hello"); \}}, to make the language easier
for beginners.

\section{The language in brief}

Java's syntax is C syntax. Primitive types are \code{int} (32 bits), \code{long} (64
bits), \code{double}, \code{float}, \code{boolean}, \code{char}; their sizes are the same
on every platform. Everything else is an \term{object}, accessed through a
\term{reference}.
\begin{lstlisting}[language=Java]
int n = 40;
double rho = 998.0;
var mu = 1.0e-3;                         // type inferred: double
String fluid = "water";                  // String is a class
double[] T = {293.1, 305.6, 299.4};      // array, fixed size
T.length;                                // 3 (a field, not a method)

import java.util.*;
List<Double> loads = new ArrayList<>(List.of(12.6, 14.0));   // growable list
loads.add(13.2);
Map<String, Double> density = new HashMap<>();
density.put("steel", 7850.0);
double r = density.get("steel");

for (double x : T) {                     // enhanced for loop
    System.out.println(x - 273.15);
}
\end{lstlisting}
Collections such as \code{List} and \code{Map} hold objects, not primitives, so a list of
numbers is a \code{List<Double>} (with a capital D, the \term{wrapper class}), and Java
converts between \code{double} and \code{Double} automatically (``autoboxing''). Strings
are compared with \code{s.equals(t)}, \emph{not} with \code{==}, which compares references.

Mathematical functions are static methods of \code{Math}: \code{Math.sqrt},
\code{Math.pow}, \code{Math.exp}, \code{Math.log}, \code{Math.sin}, \code{Math.PI}.
Formatted output uses \code{System.out.printf} with C-style format specifiers and \code{\%n}
for a platform-independent newline. Errors are reported with exceptions; Java
distinguishes \emph{checked} exceptions (such as \code{IOException}), which a method must
either catch or declare with \code{throws}, from unchecked ones (such as
\code{ArithmeticException}).

\section{Case study: the falling skydiver}

\codefile{Java}{java/Drag.java}{Euler solution of the drag equation}
\outfile{java/Drag.out}
This is almost identical to the C\# version. Constants are declared \code{static final}
and, by convention, written in upper case. Java arrays are automatically filled with zeros
when created.

\section{Object-oriented design with interfaces}

An \term{interface} lists methods that a class promises to provide, without saying how.
Any class that \code{implements} the interface can be used wherever the interface is
expected. The example below computes the heat loss through a composite wall, a classic
heat-transfer problem. Heat flows through a series of thermal resistances: convection at
the inner surface, conduction through each layer ($R = L/(kA)$), and convection at the
outer surface ($R = 1/(hA)$). The total resistance is the sum, and the heat flow is
$\dot Q = \Delta T / R_\mathrm{total}$.
\codefile{Java}{java/Thermal.java}{interfaces and records: a composite wall}
\outfile{java/Thermal.out}
The \code{record} keyword (Java 16+) defines a class whose fields are fixed at creation,
and generates the constructor and accessor methods automatically. Both record types
implement \code{ThermalResistance}, so the loop treats them uniformly. The output shows a
typical engineering insight: the 10~cm insulation layer contributes 84\% of the total
resistance, so it dominates the heat loss.

\section{Streams}

Java's \term{Stream API} (Java 8+) processes collections as pipelines of operations,
similar to C\#'s LINQ and to chained pandas operations:
\codefile{Java}{java/Streams.java}{statistics with streams}
\outfile{java/Streams.out}
The notation \code{Double::parseDouble} is a \term{method reference}, a short way of
writing the lambda \code{s -> Double.parseDouble(s)}.

\section{The JVM family and engineering uses}

The JVM runs several other languages that can use all Java libraries:
\textbf{Kotlin} (now Google's preferred language for Android, with a more concise
syntax), \textbf{Scala} (functional and object-oriented, used in big data through Apache
Spark), \textbf{Groovy} (scripting; used by the Gradle build tool) and \textbf{Clojure} (a
Lisp).

In engineering contexts you may meet Java in:
\begin{itemize}
\item \textbf{image analysis}: ImageJ and Fiji, widely used for microscopy and materials
  characterisation, are written in Java and extended with Java plug-ins (or macros, \py and
  other scripting languages);
\item \textbf{data pipelines}: Hadoop, Spark and Kafka process the large data streams of
  sensor networks and digital twins;
\item \textbf{simulation tools}: agent-based and traffic simulators (MATSim, AnyLogic),
  discrete-event simulation, and some optimisation libraries;
\item \textbf{mobile apps} for field data collection on Android;
\item \textbf{industrial software}: many enterprise systems (maintenance management,
  manufacturing execution systems) have Java back ends.
\end{itemize}
For numerical work, libraries such as Apache Commons Math, EJML and ojAlgo provide linear
algebra and statistics, but the ecosystem is much smaller than \py's.

\section{Strengths and weaknesses}

\begin{center}
\small
\begin{tabularx}{\linewidth}{@{}>{\raggedright\arraybackslash}X>{\raggedright\arraybackslash}X@{}}
\toprule
\textbf{Strengths} & \textbf{Weaknesses}\\
\midrule
Portable bytecode: runs anywhere with a JVM & Verbose; much ``boilerplate'' code\\
Memory safe, garbage collected & JVM start-up time and memory use\\
Fast after JIT warm-up & No operator overloading: matrix code is clumsy (\code{a.add(b.multiply(c))})\\
Huge, stable ecosystem; excellent tools & Small scientific-computing community\\
Very large job market & Not suited to embedded real-time control\\
\bottomrule
\end{tabularx}
\end{center}

\begin{pitfall}[Common pitfall: Java is not JavaScript]
Despite the name, Java and JavaScript (next chapter) are completely different languages.
JavaScript was given its name in 1995 for marketing reasons, when Java was fashionable.
As programmers like to say, ``Java is to JavaScript as car is to carpet''.
\end{pitfall}

\section{Further reading}
\begin{itemize}
\item The official Java tutorials, \url{https://dev.java/learn/}.
\item C.~Horstmann, \emph{Core Java}, 13th ed., Oracle Press, 2024.
\item J.~Bloch, \emph{Effective Java}, 3rd ed., Addison-Wesley, 2018.
\item R.~Sedgewick and K.~Wayne, \emph{Computer Science: An Interdisciplinary Approach},
  Addison-Wesley, 2016 --- a first course that uses Java and scientific examples.
\end{itemize}

\begin{exercises}
\item Run \file{Drag.java} with \code{java Drag.java}. Then add a command-line argument for
  the time step and print the maximum error for $\Delta t = 0.5$, $0.25$ and $0.125$.
\item In \file{Thermal.java}, how thick must the insulation be to reduce the heat loss to
  50~W? Find the answer by trial in the program, then check it by hand.
\item Add a \code{Radiation} record to \file{Thermal.java} with a linearised radiative
  resistance $R = 1/(h_r A)$, where $h_r = 4\varepsilon\sigma T_m^3$. Use the
  Stefan--Boltzmann constant $\sigma = 5.67\times10^{-8}\unit{W/(m^2\,K^4)}$ and an
  emissivity of $\varepsilon = 0.9$, a typical value for building surfaces. The mean
  temperature $T_m$ must be in kelvin, not in degrees Celsius, before you cube it.
\enlargethispage{3\baselineskip}
\item Explain the difference between \code{int} and \code{Integer} in Java. Why can you
  not write \code{List<int>}?
\item Write a Java program that prints the multiplication table up to $10\times10$, nicely
  aligned with \code{printf}.
\end{exercises}
